% !TeX root = ../main.tex
%------------------------------------------------------------------------------
%  Linings, Waterproofing and Durability
%  Sources: Muir Wood (1975); Duddeck & Erdmann (1985); Einstein & Schwartz
%  (1979); Janssen joint model; ITA WG2 segmental lining guidelines (2019);
%  fib and ACI guidance on fibre-reinforced segments.
%------------------------------------------------------------------------------
\chapter{Linings, Waterproofing and Durability}
\label{ch:linings}

The lining is the component of an underground structure required to remain
serviceable for the whole design life, commonly a century, in a wet and
sometimes chemically aggressive environment to which access is limited. This
chapter treats two lining families. A segmental ring is erected within a
shield and serves simultaneously as the reaction frame for the machine, as
primary support and as the permanent structure, three functions whose design
drivers conflict. A cast in-situ lining is placed behind a primary support
that has already accommodated the initial ground response, so the governing
question becomes how much of the long-term load it must carry alone.
Waterproofing and durability are considered together in
Section~\ref{sec:waterproofing-durability}, with fire and fit-out, since both
follow from the hydrogeological conditions established in
Chapter~\ref{ch:hydrogeology} and, more than any structural check, determine
whether the lining reaches its design life.

\section{Segmental ring design}
\label{sec:segmental-ring}

\subsection{Ring function, geometry and materials}
\label{sec:ring-function}

The segmental ring is asked to do three incompatible things. As the machine's
reaction frame it must carry concentrated thrust from the jacks, which loads it
axially and locally at the jack shoes. As primary support it must accept ground
load from the moment it leaves the shield tail, before the annulus grout has
gained strength. As the permanent structure it must carry ground and water load
for the design life and remain watertight.

Geometry follows from erection rather than from statics. A ring is divided into
segments small enough to handle and erect inside the shield, with a key segment
inserted last (Figures~\ref{fig:ring-build} and~\ref{fig:r3d-ring-build}). Universal rings, tapered on both faces, negotiate curves and
correct alignment drift by rotating the build position from ring to ring; the
taper must be sufficient for the tightest design radius plus an allowance for
correcting accumulated error. Segment count is a trade: more segments mean
easier handling and more joints, and joints are where stiffness is lost and
water enters~\citep{ita2019segmental}.

\begin{figure}[htbp]
\centering
\datafigure[\textwidth]{ch19/ring-build}
\caption{A ring of five $67.5^\circ$ segments and a $22.5^\circ$ key (K), and the developed view of four successive rings, each rotated at build so that no radial joint lies within $22.5^\circ$ of one in the neighbouring ring (schematic).}
\label{fig:ring-build}
\end{figure}

\begin{figure}[htbp]
\centering
\renderfigure{ring_build}{%
  \rlab[south]{0.27,0.80}{0.14,0.93}{Built rings}
  \rlab[west]{0.66,0.78}{0.78,0.80}{Next ring}
  \rlab[west]{0.53,0.43}{0.78,0.40}{Key}
}
\caption{Rings 1 to 3 of Figure~\ref{fig:ring-build} ($R = 3.4$~m, $t = 0.3$~m, $b = 1.5$~m). Rings 1 and 2 are built, their radial joints staggered by the rotation between builds; ring 3 is pulled radially apart into its five $67.5^\circ$ segments (malachite shades) and the $22.5^\circ$ key (carmine), which goes in last. Separations and joint widths schematic.}
\label{fig:r3d-ring-build}
\end{figure}

Materials have shifted markedly. Conventionally reinforced concrete remains
standard for large diameters and high loads. Steel fibre reinforced concrete
has displaced it for many smaller rings, because the governing actions are
often the transient ones (demoulding, handling, storage, jack thrust),
which produce local rather than flexural demands that fibres address well, and
because eliminating cage fabrication removes cost, congestion and a corrosion
path. Hybrid arrangements use fibres for the general section and bars only
where bursting demands them~\citep{ita2019segmental}.

\subsection{Load cases and ground loading}
\label{sec:segment-loads}

The permanent ground and water load is the case designers concentrate on and
the one that cracks fewest segments. Most segment damage originates in the
transient cases: demoulding at low concrete strength, handling and stacking,
erection inside the shield, and thrust from the jacks bearing on the
just-completed ring~\citep{ita2019segmental}. Grout pressure during annulus filling is a further
transient case, and it acts on a ring that does not yet have any ground
restraint. Figure~\ref{fig:segment-actions} gives the actions in a segment
stacked or lifted on supports at its ends.

\begin{figure}[htbp]
\centering
\datafigure[\textwidth]{ch19/segment-actions}
\caption{Axial force, shear and moment in a segment on a pin and a roller at its ends under self-weight, a stacking or handling case: radius 3.4~m to the centroid, $67.5^\circ$, 1.5~m wide, 0.3~m thick, 25~kN/m\textsuperscript{3}. The crown moment of 20.6~kNm gives 0.92~MPa of tension at the intrados, to be set against the tensile strength at that age.}
\label{fig:segment-actions}
\end{figure}

Ground loading itself requires a decision about arching. Full overburden
pressure is conservative for the maximum-load case and, at depth, grossly so:
the ground arches over the opening and carries most of its own weight.
Minimum ground load with full water pressure is a separate governing case, so
the load schedule has to bracket both extremes and not maximise one. Terzaghi's arching
reduction and its descendants estimate the reduced load, and the choice between
them is not a detail: the ratio between full and arched overburden at moderate
depth can exceed three. An arching-reduced load is the input to a load-based
or beam--spring check. Continuum and numerical interaction analyses derive
the load themselves, from the in-situ stress and the stiffness contrast, and
the two routes are not combined by applying a reduced load to an interaction
model.

Water load must be taken at the recovered water table rather than the
construction-period level (Chapter~\ref{ch:hydrogeology}), and for a tanked
ring it acts inward as pressure and, in shallow ground, upward as buoyancy on
an empty tunnel. Longer-term additions (future surcharge, adjacent
construction, seismic actions in Chapter~\ref{ch:seismic}) belong in the
load schedule from the outset, because strengthening a ring after
construction means building a secondary lining inside it, at the cost of the
internal envelope.

\subsection{Analysis methods}
\label{sec:lining-analysis}

Three families of method are in use, and each answers a slightly different
question.

\emph{Continuum solutions} treat the lining as an elastic ring embedded in an
elastic medium and solve the interaction in closed form. The solution of \citet{muirwood1975},
its development by \citet{duddeck1985structural}, and the formulation of
\citet{einstein1979simplified}
all express the result through two dimensionless ratios: a
compressibility ratio comparing the ground's stiffness to the lining's axial
stiffness, and a flexibility ratio comparing it to the lining's bending
stiffness (Section~\ref{sec:charts-and-lining-solutions},
Figure~\ref{fig:lining-ratios}). Those ratios carry
the central insight of lining design. When the
lining is flexible relative to the ground it deforms with the ground and
carries thrust with little moment; when it is stiff it resists the ground and
attracts large moments. A flexible lining therefore attracts less moment, but
it also converges more. Where a limit on distortion or on surface movement
governs, the advantage reverses, and in soft ground a ring that is too thin
for its radius is governed by buckling and not by bending. Within those
limits the flexible lining is the better structure, and the instinct to
stiffen is usually wrong.

\begin{figure}[htbp]
\centering
\datafigure[\textwidth]{ch19/lining-ratios}
\caption{Full-slip coefficients of \citet{einstein1979simplified} (Appendix~\ref{app:closed-form}): the moment ratio against $F$, and the thrust ratio for $k = 1$ against $C$ for a flexible lining. Markers: the 0.5~m lining of Chapter~\ref{ch:ccm} in GSI~60 ground (circle, $F \approx 3400$, $C \approx 3.2$) and GSI~30 ground (square, $F \approx 540$, $C \approx 0.50$); both carry thrust with almost no moment.}
\label{fig:lining-ratios}
\end{figure}

The interface assumption matters. Full-slip and no-slip solutions bracket the
real condition, and moments differ substantially between them; a grouted
annulus is closer to no-slip than designers often assume.

\emph{Bedded beam-spring models} discretise the ring into beam elements
supported by radial springs representing the ground, usually acting in
compression only so the ring can separate from the ground where it deflects
inward. They accommodate non-uniform loading and joint stiffness that
closed-form solutions cannot, at the cost of a strong sensitivity to the
assumed spring stiffness, which is the least reliable input.

\emph{Numerical models} embed the lining in a continuum
(Chapter~\ref{ch:numerical-modelling}) and represent staged installation. They
are necessary where ground and lining interact in ways the simpler methods
cannot capture: twin tunnels, non-uniform ground, staged excavation. They are
also where a spurious precision is easiest to generate.

The reasonable practice is to design with a closed-form solution
(Section~\ref{sec:charts-and-lining-solutions}), understand
the result through the compressibility and flexibility ratios, and use a
numerical model to check the cases the closed-form solution cannot address.

\section{Joints, connections and build quality}
\label{sec:joints-and-build}

\subsection{Joint behaviour}
\label{sec:joint-behaviour}

Radial joints are the reason a segmental ring is not a monolithic ring.
Contact is over a reduced width, not the full section, so the joint transmits
thrust with limited rotational stiffness. The model of \citet{janssen1983gelenk} treats the joint as a
reduced-height section in contact, giving a rotational stiffness that decreases
as the eccentricity of the thrust increases: as the joint opens on one side,
the contact width reduces, stiffness falls, and moment sheds to adjacent
sections (Figure~\ref{fig:joint-stiffness}).

\begin{figure}[htbp]
\centering
\datafigure[\textwidth]{ch19/joint-stiffness}
\caption{A radial joint carrying thrust $N$ through a contact band of height $h_j$ in a segment of thickness $t_c$ (schematic), and the rotational stiffness of the model of \citet{janssen1983gelenk} relative to its closed value $k_{\varphi 0} = E b h_j^2/12$: constant until the joint opens at $e = h_j/6$, then $13.5\,(e/h_j)(1 - 2e/h_j)^2$ times it.}
\label{fig:joint-stiffness}
\end{figure}

That behaviour is beneficial globally and costly locally
(Section~\ref{sec:local-effects-build}). A ring with articulating joints redistributes
moment away from its most loaded points and behaves more flexibly than its
section properties suggest, which, by the stiffness argument of
Section~\ref{sec:lining-analysis}, reduces the load it attracts. Muir Wood's effective inertia correction captures this approximately
for hand calculation; a beam-spring or numerical model can represent it
directly through rotational springs.

Circumferential joints couple successive rings. Where staggered so that radial
joints do not align, the coupling lets a heavily loaded ring share load with
its neighbours, and the ring effectively acts as part of a tube rather than in
isolation. This is a real strengthening effect and is routinely neglected;
where it is relied upon, the coupling must be detailed to transfer shear.

\subsection{Connections and watertightness}
\label{sec:connections-watertightness}

Connections serve during erection and, sometimes, permanently. Their essential
function is to hold segments in position and compress the gasket while the ring
is built and before the annulus grout sets. Bolted connections do this
reliably and leave pockets that must be sealed and that interrupt the section.
Dowel and cam systems engage automatically as the ring is pushed together,
which is faster and eliminates the pockets, and the trend is firmly towards
bolt-free systems.

Watertightness depends on the gasket, and the gasket depends on the build.
EPDM profiles seal by compression, and their performance envelope is expressed
as the water pressure sealed for a given combination of gap and offset between
adjacent segments. Both gap and offset are build-quality outcomes. A gasket
adequate for the design head at perfect build may leak at the same head with a
few millimetres of offset, and hydrophilic gaskets, which swell on wetting,
are used where tolerances cannot be guaranteed, at the cost of requiring
water to activate.

The design consequence is that the watertightness specification and the build
tolerance specification are one requirement, not two, and a gasket selected
without an achievable tolerance is a leak already designed in.

\subsection{Local effects and build quality}
\label{sec:local-effects-build}

Concentrated loads govern more segment reinforcement than global actions do.
Jack shoes bear on a small area of the ring face, and radial joints transmit
thrust over a reduced contact width; both generate transverse tension behind
the loaded area that will split the segment unless it is reinforced against
bursting. Classical strut-and-tie and partially loaded area methods size that
reinforcement, and it is frequently the governing reinforcement in an otherwise
lightly loaded ring.

Build defects have names because they recur. \emph{Ovalisation} is departure
from circularity, which introduces moments the design did not consider.
\emph{Lipping} is a radial step between adjacent segments; \emph{stepping} is
the same between rings. \emph{Birdsmouthing} is a joint open on one face,
concentrating thrust on a reduced width and inviting spalling. Each degrades
both structural performance and gasket seating.

Damage and repair complete the sequence. Cracking has diagnostic value:
longitudinal cracks near a joint usually indicate bursting or eccentric
thrust; radial cracks at mid-segment indicate flexure; corner spalling
indicates lipping or point bearing. Acceptance criteria should be set against
durability and watertightness rather than aesthetics, and repair (injection,
patching, local strengthening) should address the cause. A ring repaired
without correcting the build practice that damaged it will produce more
damaged rings.

\section{Permanent linings}
\label{sec:permanent-linings}

\subsection{Lining philosophies}
\label{sec:lining-philosophies}

Three philosophies divide the field, and the choice determines the verification
burden more than the structural quantities.

A \emph{double-shell} lining treats the primary support as sacrificial: the
secondary lining is designed for the full long-term load as though the primary
contributed nothing. It is the conservative and traditional approach, it wastes
the primary lining's capacity, and it requires the least evidence.

A \emph{composite} lining assumes primary and secondary act together. It is
more economical and lower in carbon, and it requires the interface to transfer
shear, the primary lining's long-term durability to be demonstrated, and the
load sharing between them to be justified.

A \emph{single-shell} lining is the permanent structure in its own right, with
no secondary lining at all. It is the most economical, and it places the entire
burden on the quality of a sprayed or segmental lining installed under
construction conditions.

The progression is one of increasing economy and increasing verification. The
error to avoid is claiming the economy without accepting the verification.

\subsection{Cast in-situ linings and composite action}
\label{sec:cast-insitu-composite}

Cast in-situ linings are placed with a travelling formwork in bays of ten to
fifteen metres, and the construction sequence shapes the design as much as the
loading does. Bay length, pour rate, formwork stripping time and the treatment
of construction joints all follow from the traveller cycle, and construction
joints are the locations where leakage and cracking concentrate.

Unreinforced linings remain viable and are common in good rock. A lining
carrying predominantly thrust with small eccentricity needs no reinforcement,
and omitting it removes a corrosion mechanism and substantial cost. What
unreinforced linings tolerate poorly is any action producing significant
moment: non-uniform ground load, seismic racking, or thermal restraint.
Reinforcement is therefore a decision about the moment regime, not a default.

Composite action requires assumptions about the interface and about time. A
bonded interface transfers shear and lets the two shells act as one thicker
section, with substantially higher bending capacity; a debonded interface lets
them act as two independent shells. Load sharing between them further depends
on their relative stiffness at the time of loading, and the primary lining's
stiffness rises with curing while its creep relaxes load into the secondary.
These are assumptions rather than measurements, and a composite design should
state them and demonstrate that the outcome is not sensitive to plausible
variation.

\subsection{Drained and tanked systems}
\label{sec:drained-tanked-systems}

The hydrogeological decision of Chapter~\ref{ch:hydrogeology} arrives here as a
detailing problem (Figure~\ref{fig:drained-tanked}). A \emph{drained} system provides a drainage layer behind the
lining and invert drains to convey water to a sump; the lining carries little
water pressure, but the system requires permanent pumping and drains that must
remain functional and accessible for the asset's life. A \emph{tanked} system
resists the full head with a continuous waterproofing barrier; there is nothing
to pump, and the lining is heavier, must be checked for uplift when empty, and
depends on a membrane that has to be installed correctly the first time.

\begin{figure}[htbp]
\centering
\datafigure[0.62\textwidth]{ch19/drained-tanked}
\caption{Drained and tanked linings below a water table (schematic). The drained half leads water through a drainage layer to an invert drain and carries little water pressure; the tanked half resists the full head through a continuous membrane, the pressure rising with depth.}
\label{fig:drained-tanked}
\end{figure}

Partial systems (drained invert with tanked arch, or drainage above a
defined level) are common compromises and require the transition to be
detailed explicitly, since that is where both systems are at their weakest.

\section{Waterproofing and durability}
\label{sec:waterproofing-durability}

\subsection{Waterproofing systems and drainage}
\label{sec:waterproofing-systems}

Sheet membranes (PVC or TPO, typically 2 to 3\,mm) are the established
tanking system, installed against a protective geotextile on a prepared
substrate and welded at overlaps. Their reliability rests entirely on the
welding, which is why double-track welds with a testable air channel between
them are standard: every metre of weld can be pressure tested and the result
recorded.

Compartmentalisation is the feature that makes a membrane maintainable.
Waterbars divide the waterproofed volume into compartments, each with grout
injection and venting tubes, so that a leak can be located to one compartment
and treated by re-injection rather than by exploratory demolition. Membranes
installed without compartmentalisation are cheaper and effectively
irreparable.

Sprayed membranes bond to both substrate and lining, which allows genuine
composite shell construction and eliminates the sheet membrane's slip plane.
They are attractive, but their long-term performance record is shorter than that
of sheet systems, with performance sensitive to substrate condition, curing and
application thickness. They should be specified with that asymmetry
acknowledged rather than treated as equivalent.

Drainage systems fail by clogging, not by structural inadequacy. Groundwater
carrying dissolved calcium precipitates calcite in the drains, and in some
geologies the deposition rate is fast enough to block a drain within a few
years. Sediment adds to it. A drainage system must therefore be designed for
cleaning (adequate diameter, rodding and flushing access at sensible
intervals, and inspectable outlets) and its maintenance requirement written
into the operating manual. A drain that cannot be cleaned is a temporary work
with a permanent function.

\subsection{Durability and serviceability}
\label{sec:durability-serviceability}

Concrete durability in tunnels is governed by groundwater chemistry
(Chapter~\ref{ch:hydrogeology}). Sulfate
attack expands and disrupts the cement matrix; acidic and carbonate-aggressive
water dissolves it; chlorides depassivate reinforcement and initiate pitting
corrosion; leaching removes calcium hydroxide and increases permeability over
time. The groundwater analysis assigns the exposure class, in a framework
such as that of EN~206~\citep{en206}, whose chemical-attack classes are
defined by groundwater composition. The class sets limits on binder type,
minimum cement content and maximum water-binder ratio, and the structural
design code sets the minimum cover for the same class. These are
specification decisions taken from investigation data rather than structural
calculation.

Cover sets the time taken for the carbonation or chloride front to reach the
reinforcement. Because penetration grows roughly with the square root of
time (Figure~\ref{fig:carbonation-cover}), a shortfall in cover shortens that time more than in proportion, and
once the concrete is cast it cannot be made good. Cover is also the parameter
most often compromised in construction. Its verification is a quality control matter that
belongs in the same category as segment build tolerance.

\begin{figure}[htbp]
\centering
\datafigure[0.62\textwidth]{ch19/carbonation-cover}
\caption{A carbonation or chloride front advancing as $x = K\sqrt{t}$, normalised by the nominal cover $c_{\mathrm{nom}}$ and the time $t_{\mathrm{nom}}$ the front takes to reach it. A cover 20\,\% short is reached in 64\,\% of the time, since the time scales with the square of the cover.}
\label{fig:carbonation-cover}
\end{figure}

Cracking controls serviceability. Crack width limits follow from watertightness
and from corrosion protection rather than from strength, and typical limits are
a fraction of a millimetre where water pressure acts. Two mechanisms dominate
in linings. Early thermal cracking arises in thick pours as hydration heat
generates expansion that is then restrained on cooling, and it is controlled by
limiting pour thickness and binder heat, or by accepting and detailing the
crack pattern. Restrained shrinkage produces the same effect more slowly. Both
are construction-sequence problems, and both are cheaper to prevent than to
inject.

\subsection{Fire resistance and design life}
\label{sec:fire-design-life}

Tunnel fires reach temperatures and heating rates far beyond building fires
(a heavy goods vehicle fire develops far more severely than the standard
cellulosic curve), and design fire curves such as the hydrocarbon curve and
the Rijkswaterstaat (RWS) curve reflect this.

The governing mechanism is explosive spalling. Rapid heating vaporises pore
water faster than it can escape through a low-permeability concrete, and the
resulting pore pressure spalls the surface in layers, exposing reinforcement
and reducing section progressively during the fire. Dense durable concrete is
the most vulnerable, so the durability specification and the fire
specification pull in opposite directions. Polypropylene micro-fibres are the
standard mitigation: they melt at low temperature and leave a capillary network
through which vapour escapes. Sacrificial boards are used where the
requirement is severe or the concrete cannot be modified.

The performance requirement must state what is required during the fire and
what is required after it. Stability during the fire for the evacuation and
intervention period is one criterion; retaining sufficient capacity to remain
in service, or to be repaired rather than replaced, is a different and often
governing one. For a tanked lining, a fire that breaches the waterproofing has
consequences that outlast the fire.

Design life ties the chapter together. A hundred to a hundred and twenty years
is now normal, and it is achieved through cover, binder selection, crack
control, waterproofing and a maintenance regime, not through strength. The
maintenance assumption is part of the design: a lining whose durability
depends on drains being flushed every five years has been designed on a promise
about the operator's behaviour, and that promise should be recorded in the
handover documentation rather than left implicit
(Chapter~\ref{ch:assets}).

\section{Internal fit-out and secondary structures}
\label{sec:fitout}

The lining is not the last structure built. Road pavements, track slabs,
walkways and cable routes, cross-passage doors, ventilation structures,
drainage channels and services all follow, and each interfaces with the lining
mechanically and geometrically.

Three interface issues recur. Fixings into the lining penetrate the cover and,
in a tanked structure, potentially the waterproofing, so the fixing
philosophy must be established before the membrane is designed, and cast-in
sockets are strongly preferable to post-drilled anchors. Secondary structures
add load, and a track slab or a road pavement is a substantial one that must
appear in the lining's load schedule rather than being added after it is fixed.
And the fit-out consumes the space that inspection and maintenance need: a
lining that cannot be inspected because a walkway was built against it has lost
a capability that Chapter~\ref{ch:assets} assumes it has.

Fit-out is also where clearance envelopes bind. The internal diameter required
is the sum of the operational envelope, the fit-out, the construction
tolerances and an allowance for convergence and future relining (of the order
quoted in Section~\ref{sec:repair-strengthening}). Since
diameter drives lining thickness, excavation volume, spoil and carbon
(Chapter~\ref{ch:sustainability}), the fit-out layout is a first-order
influence on the whole project rather than a detail settled late.

\chapterreferences
